\section{Big Sleep 案例复盘：从假设到崩溃证据的完整闭环}

\subsection{项目目标与任务定义}
Big Sleep（课程中也关联 Project Naptime）聚焦一个可明确验证的目标：给定代码库，生成输入触发 sanitizer crash。这个目标把 ``发现漏洞'' 转化为可执行搜索问题，既能自动评估，也能保留工程真实性。

\begin{figure}[H]
\centering
\includegraphics[width=0.88\textwidth]{frame-06-big-sleep.jpg}
\caption{Big Sleep 项目介绍画面：将漏洞发现转为可执行触发任务 \protect\footnotemark}
\end{figure}
\footnotetext{视频画面时间区间：01:14:58--01:15:35。}

\begin{importantbox}{Big Sleep 的关键建模}
目标不是让模型口头解释 ``这里有漏洞''，而是让系统产出一个可复现输入，使程序在 sanitizer/调试器下产生客观异常。换言之，结论必须被执行环境认可。
\end{importantbox}

\subsection{三类核心工具}
课程展示了该系统的三类工具：
\begin{itemize}
    \item 代码浏览工具：按符号检索、跳转定义、跟踪引用；
    \item 输入生成工具：让 Agent 写 Python 程序构造复杂输入，而非手写长字节串；
    \item 调试执行工具：运行目标程序并返回 watchpoint/breakpoint 对应表达式值。
\end{itemize}

\begin{knowledgebox}{为什么让 Agent 生成 Python 输入程序}
在安全场景，输入常包含二进制结构、嵌套格式或校验字段。让 Agent 写 ``生成器程序'' 比直接输出原始字节更稳健，也更利于后续迭代修改。
\end{knowledgebox}

\begin{warningbox}{把调试器当交互式聊天窗口并不高效}
课堂中的工具设计并非每步停住人工式单步调试，而是 ``整程序运行 + 回传命中的观察点值''。这减少了交互开销，同时保留关键证据，更适合自动化循环。
\end{warningbox}

\subsection{轨迹式问题求解}
课程演示的轨迹非常典型：先做代码浏览并形成漏洞假设，然后输出漏洞报告与利用思路；接着生成输入并运行；若未崩溃则归因失败原因、修正输入、再次运行，直到触发异常。这个过程体现了 Agent 的核心优势：\textbf{失败不是终点，而是下一轮搜索的监督信号}。

\begin{figure}[H]
\centering
\includegraphics[width=0.88\textwidth]{frame-08-wrapup.jpg}
\caption{讲座收尾对安全 Agent 现状与前景的总结 \protect\footnotemark}
\end{figure}
\footnotetext{视频画面时间区间：01:25:45--01:26:20。}

\begin{importantbox}{结果层面的关键事实}
课程披露了一个真实世界结果：在 SQLite 中发现并披露了真实漏洞。更重要的是，项目还展示了 variant analysis 能力，即基于已知漏洞模式在同项目内寻找相似变体。
\end{importantbox}

\subsection{与传统 fuzzing 的比较}
主讲人提到一个很有说服力的细节：在知道漏洞位置后，再用通用 fuzzer 反向尝试，长时间运行（课程中提到约 150 小时量级）仍很难直接找到同一漏洞。这个对比说明 Agent 语义引导在深层路径探索中可能提供增益。

\begin{knowledgebox}{Agentic 搜索的潜在价值}
\begin{itemize}
    \item 可利用先验知识（历史漏洞、补丁语义）做定向搜索；
    \item 可在失败后重写输入生成策略；
    \item 可结合调试反馈缩小可疑路径空间。
\end{itemize}
\end{knowledgebox}

\begin{warningbox}{不要把单案例成功外推为“全面超越”}
Big Sleep 证明了可行性，但不意味着所有项目、所有漏洞类型都已可自动化。讲座结论是 ``刚起步、空间很大''，而不是 ``问题已解决''。
\end{warningbox}

\subsection{本章小结}
Big Sleep 的价值在于示范了完整闭环：目标可执行、工具可调用、轨迹可迭代、结果可验证。它把安全 Agent 从概念验证推进到可复现工程实践。

